\documentclass[11pt,a4paper]{article}
\usepackage[T1]{fontenc}
\usepackage{lmodern}
\usepackage[margin=27mm]{geometry}
\usepackage{amsmath,amssymb,amsthm,mathtools}
\usepackage{booktabs,array,longtable}
\usepackage{microtype}
\usepackage{xcolor}
\usepackage{hyperref}
\usepackage{fancyhdr}
\hypersetup{colorlinks=true,linkcolor=blue!45!black,citecolor=blue!45!black,
 urlcolor=blue!45!black,pdftitle={Counting terms in derivatives of x to the x: OEIS A293239},
 pdfauthor={OpenAI ChatGPT},pdfsubject={Exact reduction, partial proofs, and computational certificates}}
\pagestyle{fancy}
\fancyhf{}
\fancyhead[L]{\small Counting terms in derivatives of $x^x$}
\fancyhead[R]{\small OEIS A293239}
\fancyfoot[C]{\thepage}
\setlength{\headheight}{14pt}
\setlength{\emergencystretch}{2em}
\numberwithin{equation}{section}
\newtheorem{theorem}{Theorem}[section]
\newtheorem{proposition}[theorem]{Proposition}
\newtheorem{lemma}[theorem]{Lemma}
\newtheorem{corollary}[theorem]{Corollary}
\theoremstyle{definition}
\newtheorem{definition}[theorem]{Definition}
\newtheorem{conjecture}[theorem]{Conjecture}
\theoremstyle{remark}
\newtheorem{remark}[theorem]{Remark}
\newcommand{\coef}[2]{[#1^{#2}]}
\newcommand{\fall}[2]{#1^{\underline{#2}}}
\newcommand{\ind}[1]{\mathbf{1}_{\{#1\}}}
\newcommand{\Z}{\mathbb Z}
\newcommand{\Q}{\mathbb Q}
\newcommand{\R}{\mathbb R}
\newcommand{\N}{\mathbb N}
\newcommand{\ZZ}{\mathcal Z}
\newcommand{\EE}{\mathcal E}
\DeclareMathOperator{\supp}{supp}
\title{\textbf{Counting terms in derivatives of $x^x$}\\[5pt]
\large An exact reduction of OEIS A293239, nonvanishing theorems,\\
computational certificates, and a corrected recurrence}
\author{Research report prepared with ChatGPT}
\date{20 September 2026}
\begin{document}
\maketitle
\begin{abstract}
The sequence A293239 counts the nonzero terms in the fully expanded $n$th
 derivative of $x^x$. Its proposed generating function is rational, and its
 proposed eventual formula is a quadratic quasipolynomial of period four.
 We reduce the counting problem exactly to the zero set of the first-kind
 Lehmer--Comtet numbers. All cancellations predicted by the conjecture are
 proved: a symmetry family and one additional zero. The converse assertion,
 excluding every further zero, is not established here; consequently this
 report does \emph{not} claim a complete proof of the principal closed formula.
 The unconditional results include an exact defect formula, upper and lower
 bounds proving quadratic-order growth, nonvanishing on every prime-offset
 diagonal with an exact prime-adic valuation, and complete zero classifications
 for the first three columns and offsets one through sixteen. Independently
 implemented exact computations confirm the formula through $n=3000$.
 A separate conjectured recurrence in the OEIS entry is false with its printed
 range: it already fails at $n=8$. For the proposed sequence its correct
 eventual starting point is $n=14$. The accompanying code and certificates
 distinguish finite verification from statements proved for infinitely many
 indices.
\end{abstract}
\clearpage
\tableofcontents
\newpage

\section{The question and the status of the answer}

For $x>0$, put $f(x)=x^x=\exp(x\log x)$. Let $a(n)$ be the number of nonzero
summands after $f^{(n)}(x)$ is expanded into the monomials
\begin{equation}\label{eq:monomials}
 x^{x-k}(\log x)^j,\qquad k,j\in\Z_{\ge0},
\end{equation}
and coefficients of identical monomials are collected. In particular, an
unexpanded expression such as $x^x(1+\log x)^n$ is not one term for this
purpose. This is the convention in OEIS A293239, introduced by Vladimir
Reshetnikov in 2017 \cite{oeis293239}.

The entry attributes the rational generating function and closed formula to
Colin Barker and the proposed asymptotic equivalent to V\'aclav
Kote\v{s}ovec. To separate a proved object from the proposed answer, define
\begin{equation}\label{eq:q}
 q(0)=1,\qquad
 q(n)=1+\frac{n(n+1)}2-\left\lfloor\frac{n-1}{4}\right\rfloor
             -\ind{n\ge8}\quad(n\ge1).
\end{equation}
The principal conjecture is $a(n)=q(n)$. For $n\ge8$, this says simply
\begin{equation}\label{eq:eventual}
 a(n)\stackrel{?}{=}\frac{n(n+1)}2-\left\lfloor\frac{n-1}{4}\right\rfloor.
\end{equation}
The proposed generating function is
\begin{equation}\label{eq:targetgf}
 Q(z):=\sum_{n\ge0}q(n)z^n
 =\frac{1+z^2+z^3+z^6-z^8+z^9+z^{12}-z^{13}}
 {(1-z)^2(1-z^4)}.
\end{equation}
The rational function in \eqref{eq:targetgf} is indeed the generating function
of the explicitly defined sequence $q$. What remains to be proved is its
identification with the derivative count $a$.

\subsection{What is proved}
The central unconditional identity will be
\begin{equation}\label{eq:introdefect}
 a(n)=q(n)-E(n),
\end{equation}
where $E(n)$ counts a precisely defined set of additional zero coefficients.
In particular, $E(n)$ is a nonnegative, nondecreasing integer. Thus the
proposed answer is an upper bound, not a formula that can be exceeded by a
previously overlooked cancellation.

We prove, for all $n\ge0$,
\begin{equation}\label{eq:introbounds}
 n+1+\left\lfloor\frac{n^2}{4}\right\rfloor\le a(n)\le q(n).
\end{equation}
Consequently $a(n)=\Theta(n^2)$. This does not establish the sharper claim
$a(n)\sim n^2/2$; that claim is equivalent to $E(n)=o(n^2)$.
We also prove infinite nonvanishing results that substantially restrict where
additional zeros could lie. The remaining gap is stated explicitly in
Conjecture~\ref{conj:zeros}.

\subsection{A correction that is unconditional}
The OEIS entry also proposes
\begin{equation}\label{eq:recprinted}
 a(n)=2a(n-1)-a(n-2)+a(n-4)-2a(n-5)+a(n-6)\qquad(n>6).
\end{equation}
This range is incorrect. Exact low-order differentiation gives
\[
 a(2)=4,\quad a(3)=7,\quad a(4)=11,\quad a(6)=21,\quad
 a(7)=28,\quad a(8)=35,
\]
whereas the right side of \eqref{eq:recprinted} at $n=8$ is
\[
 2\cdot28-21+11-2\cdot7+4=36\ne35.
\]
For $q$, the recurrence holds for every $n\ge14$, and it fails at $n=13$.
This correction does not disprove the proposed generating function or
closed formula. Section~\ref{sec:gf} explains the distinction.

\section{A canonical polynomial for each derivative}

Set $t=x^{-1}$ and $L=\log x$. There is a polynomial $P_n\in\Z[t,L]$ such that
\begin{equation}\label{eq:Pdefinition}
 f^{(n)}(x)=x^x P_n(x^{-1},\log x).
\end{equation}
Differentiating this expression and using $t'=-t^2$, $L'=t$, and
$(x^x)'=x^x(1+L)$ gives the following exact recurrence.

\begin{proposition}\label{prop:Prec}
The derivative polynomials satisfy
\begin{equation}\label{eq:Prec}
 P_0=1,\qquad P_{n+1}=(1+L)P_n+t\frac{\partial P_n}{\partial L}
                         -t^2\frac{\partial P_n}{\partial t}.
\end{equation}
In particular, they have integer coefficients and can be computed using
integer arithmetic only.
\end{proposition}

For example,
\begin{align*}
 P_1&=1+L,\\
 P_2&=1+2L+L^2+t,\\
 P_3&=1+3L+3L^2+L^3+3t+3tL-t^2.
\end{align*}
Thus $a(3)=7$, consistently with the example in the entry.

\begin{lemma}\label{lem:independence}
The monomials in \eqref{eq:monomials} are linearly independent over $\R$
when only finitely many are used. Consequently the collected expansion is
unique, and $a(n)=|\supp P_n|$.
\end{lemma}
\begin{proof}
Divide a putative relation by the positive function $x^x$ and set $x=e^s$.
The result is $\sum_{k=0}^K e^{-ks}p_k(s)=0$, with polynomial $p_k$.
If some $p_k$ is nonzero, let $k_0$ be the smallest such index. Multiplication
by $e^{k_0s}$ shows that $p_{k_0}(s)$ tends to zero as $s\to+\infty$, because
all remaining summands decay exponentially times a polynomial. No nonzero
polynomial has that property. This is a contradiction.
\end{proof}

The independence statement matters: the variables $t$ and $L$ are related
by $t=e^{-L}$ after substitution, but this relation causes no additional
finite polynomial identities.

\section{The Lehmer--Comtet coefficient triangle}

\subsection{Definition and generating function}
Put
\begin{equation}\label{eq:phi}
 \Phi(z)=(1+z)\log(1+z)
 =z+\sum_{d\ge1}\frac{(-1)^{d-1}}{d(d+1)}z^{d+1}.
\end{equation}
Define
\begin{equation}\label{eq:bdef}
 b(m,r)=\frac{m!}{r!}\coef{z}{m}\Phi(z)^r
 \quad(m,r\ge0).
\end{equation}
Here $b(0,0)=1$, $b(m,0)=0$ for $m>0$, and $b(m,r)=0$ when $r>m$.
These are the first-kind Lehmer--Comtet numbers, tabulated in OEIS A008296
\cite{oeis8296}. The general derivative problem and these numbers belong to
the classical literature of Comtet, Lehmer, and Gould
\cite{comtet,lehmer,gould}. A recent treatment of their identities is
Jeong's preprint \cite{jeong}. We derive the identities needed here rather
than assuming a nonvanishing conclusion from those formulas.

\begin{theorem}[Exact coefficient formula]\label{thm:expansion}
For every $n\ge0$,
\begin{equation}\label{eq:expansion}
 \boxed{\displaystyle
 f^{(n)}(x)=x^x\sum_{m=0}^n\binom nm (\log x)^{n-m}
                         \sum_{r=0}^m b(m,r)x^{r-m}.}
\end{equation}
Equivalently,
\begin{equation}\label{eq:cformula}
 \coef{t}{k}\coef{L}{j}P_n(t,L)
 =\binom nj\,b(n-j,n-j-k),
\end{equation}
with out-of-range entries interpreted as zero.
\end{theorem}
\begin{proof}
Taylor expansion at a fixed $x>0$ gives
\[
 \sum_{n\ge0}\frac{f^{(n)}(x)}{x^x}\frac{u^n}{n!}
 =\frac{(x+u)^{x+u}}{x^x}
 =\exp(u\log x)\exp\bigl(x\Phi(u/x)\bigr).
\]
Using \eqref{eq:bdef}, expand the second exponential:
\begin{align*}
 \exp\bigl(x\Phi(u/x)\bigr)
 &=\sum_{r\ge0}\frac{x^r}{r!}\Phi(u/x)^r\\
 &=\sum_{m\ge0}\frac{u^m}{m!}\sum_{r=0}^m b(m,r)x^{r-m}.
\end{align*}
Multiplication by $\exp(u\log x)$ and extraction of the coefficient of
$u^n$ proves \eqref{eq:expansion}. All manipulations can also be performed
formally, since only finitely many terms affect any prescribed coefficient.
\end{proof}

\subsection{A shifted falling-factorial formula}
Let
\[
 F_m(y)=\fall{y}{m}=\prod_{h=0}^{m-1}(y-h),\qquad F_0(y)=1.
\]

\begin{proposition}\label{prop:falling}
For $0\le r\le m$,
\begin{equation}\label{eq:falling}
 b(m,r)=\frac{F_m^{(r)}(r)}{r!}
       =\coef{u}{r}\prod_{h=0}^{m-1}(u+r-h).
\end{equation}
In particular, all $b(m,r)$ are integers.
\end{proposition}
\begin{proof}
The binomial-series identity is
$(1+z)^y=\sum_{m\ge0}F_m(y)z^m/m!$.
Differentiate $r$ times in $y$ and set $y=r$. The left side becomes
$(1+z)^r\log(1+z)^r=\Phi(z)^r$. Comparing coefficients gives the first
formula. The second is the polynomial Taylor formula at $y=r$.
\end{proof}

Writing $F_m(y)=\sum_{\ell=0}^m s(m,\ell)y^\ell$ in terms of the signed
Stirling numbers of the first kind also gives
\begin{equation}\label{eq:stirling}
 b(m,r)=\sum_{\ell=r}^m\binom\ell r r^{\ell-r}s(m,\ell).
\end{equation}
This is another exact formula, but its alternating summands do not by
themselves resolve whether the result can vanish.

\subsection{An integer recurrence}
\begin{proposition}\label{prop:brec}
For $m\ge2$ and $1\le r\le m$,
\begin{equation}\label{eq:brec}
 b(m,r)=(r-m+1)b(m-1,r)+b(m-1,r-1)+(m-1)b(m-2,r-1).
\end{equation}
Together with rows $b(0,0)=1$ and $(b(1,0),b(1,1))=(0,1)$, this determines
the entire triangle.
\end{proposition}
\begin{proof}
Let $B_r(z)=\Phi(z)^r/r!$. Since
$(1+z)\Phi'(z)=1+z+\Phi(z)$, we have
\[
 (1+z)B_r'(z)=(1+z)B_{r-1}(z)+rB_r(z).
\]
Compare coefficients of $z^{m-1}/(m-1)!$. The result is
\[
 b(m,r)+(m-1)b(m-1,r)
 =b(m-1,r-1)+(m-1)b(m-2,r-1)+r b(m-1,r),
\]
which rearranges to \eqref{eq:brec}.
\end{proof}

For reference, the first rows, omitting column zero, are
\begin{center}
\setlength{\tabcolsep}{4pt}
\small
\begin{tabular}{r|rrrrrrrrr}
$m\backslash r$&1&2&3&4&5&6&7&8&9\\\hline
1&1\\
2&1&1\\
3&$-1$&3&1\\
4&2&$-1$&6&1\\
5&$-6$&0&5&10&1\\
6&24&4&$-15$&25&15&1\\
7&$-120$&$-28$&49&$-35$&70&21&1\\
8&720&188&$-196$&49&0&154&28&1\\
9&$-5040$&$-1368$&944&0&$-231$&252&294&36&1
\end{tabular}
\end{center}

\section{The exact zero-count reduction}

\begin{definition}
For $n\ge0$, let
\[
 Z(n)=\#\{(m,r):1\le r\le m\le n,\ b(m,r)=0\}.
\]
Column zero is not included: its zeros for $m>0$ are forced by the definition
and never represent potential nonzero derivative terms.
\end{definition}

\begin{theorem}\label{thm:count}
For every $n\ge0$,
\begin{equation}\label{eq:count}
 \boxed{\displaystyle a(n)=1+\frac{n(n+1)}2-Z(n).}
\end{equation}
\end{theorem}
\begin{proof}
In \eqref{eq:expansion}, the pair $(m,r)$ gives the monomial
\[
 x^{x-(m-r)}(\log x)^{n-m}
\]
with coefficient $\binom nm b(m,r)$. The map
$(m,r)\mapsto(m-r,n-m)$ is injective for fixed $n$: the second coordinate
recovers $m$, and the first then recovers $r$. Thus different pairs do not
cancel each other. The factor $\binom nm$ is positive. There is one nonzero
entry at $(0,0)$, and $m$ potentially nonzero entries $1\le r\le m$ in every
row $m\ge1$. Removing exactly the zero entries proves the formula.
\end{proof}

This argument isolates the main difficulty. Once the triangle is known,
term counting is easy; proving that its list of zeros is complete is not
a consequence of merely writing down the triangle's recurrence.

\subsection{Cancellations that are proved}
\begin{proposition}[Symmetry zeros]\label{prop:symmetry}
For every integer $h\ge1$,
\begin{equation}\label{eq:symmetry}
 b(4h+1,2h)=0.
\end{equation}
\end{proposition}
\begin{proof}
Center the product in \eqref{eq:falling} at $r=2h$:
\[
 F_{4h+1}(2h+u)=u\prod_{s=1}^{2h}(u^2-s^2).
\]
This polynomial is odd. Its coefficient of the even power $u^{2h}$ is zero.
\end{proof}

This family is not claimed as new: it appears as Proposition 3.10(iii) in
Jeong \cite{jeong}, with attribution there to Theorem 7 of Lehmer
\cite{lehmer}. Its elementary proof is included to make the reduction
self-contained.

There is one further zero required by the proposed formula:
\begin{equation}\label{eq:exception}
 b(8,5)=0.
\end{equation}
It follows immediately from the factorization proved in
Section~\ref{sec:diagonals}:
\[
 b(m,m-3)=\frac{m(m-1)(m-2)(m-3)(m-5)(m-8)}{48}.
\]
The root $m=5$ is the first symmetry zero; $m=8$ is genuinely additional.

\subsection{The precise missing assertion}
\begin{conjecture}[Complete zero classification]\label{conj:zeros}
For integers $1\le r\le m$,
\begin{equation}\label{eq:zeroconj}
 b(m,r)=0\quad\Longleftrightarrow\quad
 (m,r)=(8,5)\ \text{or}\ (m,r)=(4h+1,2h)\text{ for some }h\ge1.
\end{equation}
\end{conjecture}

The right-to-left implication is proved above. The left-to-right implication
is the unproved step in this report.

Let
\[
 \ZZ_0=\{(8,5)\}\cup\{(4h+1,2h):h\ge1\}
\]
be the proved set of zeros, and define
\begin{align*}
 e_m&=\#\{r:1\le r\le m,\ b(m,r)=0,\ (m,r)\notin\ZZ_0\},\\
 E(n)&=\sum_{m=1}^n e_m.
\end{align*}
Thus $E$ counts only zeros not in the proved list.

\begin{theorem}[Exact defect formula]\label{thm:defect}
For every $n\ge0$,
\[
 a(n)=q(n)-E(n).
\]
The following are equivalent: Conjecture~\ref{conj:zeros}; $E(n)=0$ for all
$n$; $a(n)=q(n)$ for all $n$; and the proposed eventual formula
\eqref{eq:eventual} for every $n\ge8$.
\end{theorem}
\begin{proof}
For $n\ge1$, the number of symmetry zeros with first coordinate at most $n$
is $\lfloor(n-1)/4\rfloor$. The zero $(8,5)$ contributes
$\ind{n\ge8}$. Substituting
\[
 Z(n)=\left\lfloor\frac{n-1}{4}\right\rfloor+\ind{n\ge8}+E(n)
\]
into \eqref{eq:count} proves the identity for $n\ge1$; $n=0$ is immediate.
The first three equivalences follow from the definitions. The eventual
formula also forces $E(n)=0$ for every $n\ge8$, and nonnegativity and
monotonicity then force $E(n)=0$ for smaller $n$ as well.
\end{proof}

A hypothetical new zero at row $M$ would decrease every count $a(n)$ for
$n\ge M$ relative to $q(n)$. It would not merely alter a single isolated
term of the sequence. This persistent-deficit property is a useful safeguard
when testing proposed proofs or numerical data.

\section{Diagonal polynomials and exact nonvanishing}\label{sec:diagonals}

\subsection{A polynomial for each fixed offset}
Define
\begin{equation}\label{eq:H}
 H(z)=\frac{\Phi(z)}z
 =1+\frac z2-\frac{z^2}{6}+\frac{z^3}{12}-\frac{z^4}{20}+\cdots,
 \qquad T_d(Y)=\coef{z}{d}H(z)^Y.
\end{equation}
For a fixed integer $d\ge0$, $T_d$ is a polynomial over $\Q$ in $Y$.

\begin{lemma}\label{lem:diagdegree}
For $m\ge d$,
\begin{equation}\label{eq:diag}
 b(m,m-d)=\fall m d\,T_d(m-d).
\end{equation}
For $d\ge1$, $T_d(Y)$ has degree $d$, leading coefficient $1/(2^d d!)$,
and a factor $Y$. Consequently $T_d(Y)/Y$ is a nonzero polynomial of
degree $d-1$.
\end{lemma}
\begin{proof}
Since $\Phi(z)^{m-d}=z^{m-d}H(z)^{m-d}$, definition \eqref{eq:bdef} gives
\eqref{eq:diag}. Now assume $d\ge1$ and write $h(z)=H(z)-1$. Then
\begin{equation}\label{eq:Tbinomial}
 T_d(Y)=\sum_{\ell=1}^d\binom Y\ell\coef{z}{d}h(z)^\ell.
\end{equation}
Every summand is divisible by $Y$. The only term of degree $d$ in $Y$
comes from $\ell=d$, and $\coef{z}{d}h(z)^d=2^{-d}$. This proves the
remaining assertions.
\end{proof}

The first three cases are
\begin{align}
 b(m,m-1)&=\frac{m(m-1)}2,\label{eq:d1}\\
 b(m,m-2)&=\frac{m(m-1)(m-2)(3m-13)}{24},\label{eq:d2}\\
 b(m,m-3)&=\frac{m(m-1)(m-2)(m-3)(m-5)(m-8)}{48}.
 \label{eq:d3}
\end{align}
For instance,
\[
 T_3(Y)=\frac{Y}{12}-\frac{Y(Y-1)}{12}
                  +\frac{Y(Y-1)(Y-2)}{48}
        =\frac{Y(Y-2)(Y-5)}{48}.
\]
This proves both zeros on the third offset and excludes all others on that
entire diagonal, not just up to a numerical cutoff.

\subsection{Prime offsets: an infinite family of nonvanishing results}
For a nonzero rational number $v$, write $v_p(v)$ for its exponent of the
prime $p$ in its prime factorization. A rational number is called
$p$-integral when its denominator in lowest terms is not divisible by $p$.

\begin{theorem}[Prime-offset nonvanishing]\label{thm:prime}
Let $p$ be any prime and let $m\ge p$. Then
\begin{equation}\label{eq:primevaluation}
 b(m,m-p+1)\ne0,\qquad
 v_p\bigl(b(m,m-p+1)\bigr)=v_p\left(\left\lfloor\frac mp\right\rfloor\right).
\end{equation}
In particular, an additional zero cannot occur on any offset $d=m-r$
for which $d+1$ is prime.
\end{theorem}
\begin{proof}
Put $d=p-1$ and $h_j=\coef{z}{j}(H(z)-1)$.
For $1\le j\le p-2$, the number
$h_j=(-1)^{j-1}/(j(j+1))$ is $p$-integral. Moreover,
\[
 p h_{p-1}\equiv1\pmod p.
\]
For odd $p$, this follows from $p h_{p-1}=-1/(p-1)$; for $p=2$, it is
$2h_1=1$.

Divide \eqref{eq:Tbinomial} by $Y$. Its $\ell=1$ summand is $h_{p-1}$.
For every $2\le\ell\le p-1$, a composition of $p-1$ into $\ell$ positive
parts has all parts at most $p-2$. Therefore every coefficient contributing
to $\coef{z}{p-1}h(z)^\ell$ is $p$-integral. Also,
\[
 \frac1Y\binom Y\ell
 =\frac{(Y-1)\cdots(Y-\ell+1)}{\ell!}
\]
has $p$-integral coefficients, because $\ell!$ is not divisible by $p$.
It follows, as a polynomial congruence over the $p$-integral rationals, that
\begin{equation}\label{eq:punit}
 p\,\frac{T_{p-1}(Y)}Y\equiv1\pmod p.
\end{equation}
For every positive integer $Y$, the left side is therefore nonzero and has
$p$-adic valuation zero. Hence
$v_p(T_{p-1}(Y))=v_p(Y)-1$.

Set $Y=m-p+1>0$ and use \eqref{eq:diag}. This yields
\[
 v_p\bigl(b(m,m-p+1)\bigr)=v_p(\fall m p)-1.
\]
Among the $p$ consecutive positive integers $m-p+1,\ldots,m$, exactly one
is divisible by $p$: it is $p\lfloor m/p\rfloor$. Its valuation is
$1+v_p(\lfloor m/p\rfloor)$, proving \eqref{eq:primevaluation}.
\end{proof}

This proof explains, for example, why the residual polynomial $3m-13$ in
\eqref{eq:d2} is a nonzero constant modulo $3$, and why analogous residual
polynomials for offsets $4,6,10,12,16$ are nonzero constants modulo
$5,7,11,13,17$, respectively.

\subsection{Finite certificates that cover infinitely many row indices}
For a fixed $d\ge1$, put
\[
 S_d(X)=T_d(X-d).
\]
The factor $X-d$ is forced by Lemma~\ref{lem:diagdegree}. When $d\ge3$ is
odd, the symmetry zero also forces $X-(2d-1)$; for $d=3$ there is in
addition the factor $X-8$.
Remove these factors from $S_d$, clear rational denominators, divide out
the greatest common divisor of the integer coefficients, and normalize the
leading coefficient to be positive. Call the resulting polynomial $R_d(X)$.
Thus, with a positive rational constant $c_d$,
\begin{equation}\label{eq:Rd}
 S_d(X)=c_d(X-d)
 \begin{cases}
 (X-5)(X-8)R_d(X),&d=3,\\
 (X-(2d-1))R_d(X),&d\ge5\text{ odd},\\
 R_d(X),&\text{otherwise}.
 \end{cases}
\end{equation}

\begin{lemma}[Modular certificate]\label{lem:modcert}
If $R\in\Z[X]$ satisfies $R(t)\not\equiv0\pmod p$ for every
$t\in\{0,1,\ldots,p-1\}$, then $R$ has no integer root.
\end{lemma}
\begin{proof}
An integer root, reduced modulo $p$, would be one of those residues and
would give a zero value modulo $p$.
\end{proof}

\begin{theorem}[Computer-assisted diagonal classification]\label{thm:16}
Conjecture~\ref{conj:zeros} holds for every pair $1\le r\le m$ with
$m-r\le16$.
\end{theorem}
\begin{proof}
For offset zero, $b(m,m)=1$. For $1\le d\le16$, compute $S_d$ exactly from
\eqref{eq:Tbinomial} and factor it as in \eqref{eq:Rd}. The supplied
\texttt{data/diagonal\_certificates.json} records the full rational
coefficients of $S_d$, the scalar $c_d$, all removed linear factors, the
full integer coefficients of $R_d$, and every value of $R_d(t)$ modulo
the certifying prime. All of those values are nonzero. The following
summary identifies the certificates:
\begin{center}
\begin{tabular}{rrrr@{\qquad}rrrr}
\toprule
$d$&$\deg R_d$&$p$&roots mod $p$&$d$&$\deg R_d$&$p$&roots mod $p$\\
\midrule
1&0&2&0&9&7&7&0\\
2&1&3&0&10&9&11&0\\
3&0&2&0&11&9&13&0\\
4&3&5&0&12&11&13&0\\
5&3&3&0&13&11&11&0\\
6&5&7&0&14&13&23&0\\
7&5&5&0&15&13&19&0\\
8&7&11&0&16&15&17&0\\
\bottomrule
\end{tabular}
\end{center}
The standard-library program \texttt{verify\_diagonals.py} reconstructs
every polynomial using rational arithmetic and checks the factorization
and each residue. The modular lemma excludes every integer root of every
$R_d$. For $m>d$, the removed factor $m-d$ cannot vanish; the remaining
removed factors give exactly the stated symmetry and exceptional zeros.
\end{proof}

These certificates are not extrapolations from finitely many values of
$m$: each modular argument excludes integer roots for \emph{all} $m$ on
its fixed diagonal. What is finite is the set of offsets treated.

For a small example not requiring the machine-readable file,
\begin{align*}
 R_4(X)&=15X^3-390X^2+3245X-8638\equiv2\pmod5,\\
 R_5(X)&=3X^3-103X^2+1118X-3776
                  \equiv2X^2+2X+1\pmod3.
\end{align*}
The second reduced polynomial takes the values $1,2,1$ at the three
residues modulo $3$. Hence there are no additional zeros on offset $4$,
and the only zero on offset $5$ is the symmetry zero $m=9$.

\section{The first three columns for all row indices}\label{sec:columns}

The diagonal results keep $m-r$ fixed. A complementary calculation keeps
$r$ fixed and proves nonvanishing in a different infinite region.

\begin{proposition}\label{prop:firstcolumns}
Column $r=1$ has no zero. Column $r=2$ has exactly the zero $m=5$.
Column $r=3$ has no zero.
\end{proposition}
\begin{proof}
For $r=1$, \eqref{eq:falling} gives
\[
 b(1,1)=1,\qquad b(m,1)=(-1)^m(m-2)!\quad(m\ge2).
\]
For $r=2$, put $t=m-3\ge0$. The shifted product is
\[
 F_m(2+u)=u(u+1)(u+2)\prod_{j=1}^{t}(u-j).
\]
The coefficient of $u^2$ is
\begin{equation}\label{eq:col2}
 b(m,2)=(-1)^m(m-3)!\,(2H_{m-3}-3),\qquad m\ge3,
\end{equation}
where $H_t=\sum_{j=1}^t1/j$ and $H_0=0$. The harmonic numbers strictly
increase, and $H_2=3/2$, so the only zero is $m-3=2$. The remaining
boundary entry is $b(2,2)=1$.

For $r=3$, the boundary entry is $b(3,3)=1$. Put $t=m-4\ge0$ and define
\[
 H_t^{(2)}=\sum_{j=1}^t\frac1{j^2},\qquad
 g(t)=6-11H_t+3\bigl(H_t^2-H_t^{(2)}\bigr).
\]
Extracting the coefficient of $u^3$ from
$u(u+1)(u+2)(u+3)\prod_{j=1}^t(u-j)$ gives
\begin{equation}\label{eq:col3}
 b(m,3)=(-1)^t t!\,g(t).
\end{equation}
The elementary identity
\begin{equation}\label{eq:gdiff}
 g(t+1)-g(t)=\frac{6H_t-11}{t+1}
\end{equation}
controls its entire sign pattern. We have $g(0)=6$, $g(1)=-5$, and
$H_3=11/6$. Thus $g$ decreases through $t=3$, has
$g(3)=g(4)=-49/6$, and is strictly increasing thereafter.
Finally, exact rational evaluation gives
\[
 g(19)=-\frac{4571462267}{97772875200}<0,\qquad
 g(20)=\frac{537826687}{1150269120}>0.
\]
It follows that $g(t)<0$ for $1\le t\le19$, whereas $g(t)>0$ for
$t\ge20$. There is no integer zero.
\end{proof}

The last argument illustrates a distinction from a numerical search: exact
values at two adjacent integers, together with a proved monotonicity
statement outside a finite initial segment, settle the entire column.

\section{Unconditional quadratic growth}\label{sec:bounds}

\begin{theorem}\label{thm:bounds}
For every integer $n\ge0$,
\begin{equation}\label{eq:bounds}
 n+1+\left\lfloor\frac{n^2}{4}\right\rfloor\le a(n)\le q(n).
\end{equation}
In particular,
\[
 \frac14\le\liminf_{n\to\infty}\frac{a(n)}{n^2}
 \le\limsup_{n\to\infty}\frac{a(n)}{n^2}\le\frac12,
 \qquad a(n)=\Theta(n^2).
\]
\end{theorem}
\begin{proof}
The upper bound is Theorem~\ref{thm:defect}. For the lower bound, count
nonzero entries along diagonals of the triangle in
Theorem~\ref{thm:count}. The main diagonal, including $(0,0)$, contributes
$n+1$ entries.

Fix $1\le d\le n-1$. The diagonal $m-r=d$ has $n-d$ relevant positions,
$m=d+1,\ldots,n$. By Lemma~\ref{lem:diagdegree}, its zero positions are
roots of the nonzero polynomial $T_d(m-d)/(m-d)$ of degree $d-1$.
There are therefore at most $d-1$ such positions, and at least
\[
 \max(0,n-d-(d-1))=\max(0,n-2d+1)
\]
nonzero entries on the diagonal. Summing over all $d\ge1$ gives
\[
 a(n)\ge n+1+\sum_{d\ge1}\max(0,n-2d+1)
       =n+1+\left\lfloor\frac{n^2}{4}\right\rfloor.
\]
The last equality follows by summing the odd numbers $1,3,\ldots,2s-1$
when $n=2s$, and the even numbers $2,4,\ldots,2s$ when $n=2s+1$.
The asymptotic bounds follow immediately.
\end{proof}

\begin{corollary}[A stronger explicit lower bound]\label{cor:primebound}
For every $n\ge0$,
\begin{equation}\label{eq:primebound}
 a(n)\ge n+1+\left\lfloor\frac{n^2}{4}\right\rfloor
       +\sum_{\substack{p\le n\\p\text{ prime}}}\min(p-2,n-p+1).
\end{equation}
\end{corollary}
\begin{proof}
On offset $d=p-1$, Theorem~\ref{thm:prime} guarantees all $n-d=n-p+1$
entries, rather than just $\max(0,n-2p+3)$. The improvement is
$\min(p-2,n-p+1)$. These diagonals are disjoint, so their improvements add.
\end{proof}

The prime-enhanced bound is useful without invoking any theorem about the
density of primes. It still does not close the factor-of-two gap in the
leading quadratic coefficient. In terms of the exact defect,
\begin{equation}\label{eq:asympeq}
 a(n)\sim\frac{n^2}{2}\quad\Longleftrightarrow\quad E(n)=o(n^2).
\end{equation}
Thus the proposed asymptotic equivalent is potentially a weaker target than
the complete zero classification: it would allow a sufficiently sparse set
of additional zeros.

\section{The rational function, quasipolynomial, and recurrence}\label{sec:gf}

\subsection{Deriving the proposed generating function}
There is one potential term at $(0,0)$ and $m$ potential terms in row $m$.
The cumulative potential count therefore has generating function
\[
 \sum_{n\ge0}\left(1+\frac{n(n+1)}2\right)z^n
 =\frac1{1-z}+\frac{z}{(1-z)^3}.
\]
The symmetry zeros are born at rows $5,9,13,\ldots$, and each remains
absent from every later derivative. Their cumulative generating function is
$z^5/((1-z)(1-z^4))$. The exceptional zero is born at row $8$ and
contributes $z^8/(1-z)$. Subtracting these two \emph{proved} families gives
\begin{align}
 Q(z)
 &=\frac1{1-z}+\frac{z}{(1-z)^3}
       -\frac{z^5}{(1-z)(1-z^4)}-\frac{z^8}{1-z}\notag\\
 &=\frac{1+z^2+z^3+z^6-z^8+z^9+z^{12}-z^{13}}
        {(1-z)^2(1-z^4)}.\label{eq:gfderive}
\end{align}
This explains every correction term in the proposed rational function.

\begin{theorem}[Exact generating-function correction]\label{thm:gfdefect}
With $A(z)=\sum_{n\ge0}a(n)z^n$ and
$\mathcal E(z)=\sum_{m\ge1}e_m z^m$,
\begin{equation}\label{eq:gfdefect}
 A(z)=Q(z)-\frac{\mathcal E(z)}{1-z}.
\end{equation}
Consequently the proposed rational function equals $A(z)$ if and only if
Conjecture~\ref{conj:zeros} is true.
\end{theorem}
\begin{proof}
Sum $a(n)=q(n)-\sum_{m\le n}e_m$ as a formal power series. Each $e_m$
contributes $e_mz^m/(1-z)$ to the cumulative sum.
\end{proof}

\subsection{The real form of the complex expression}
For $n\ge8$, division into the four residue classes gives
\begin{equation}\label{eq:residueq}
 q(n)=\frac{n^2}{2}+\frac n4+
 \begin{cases}
 1,&n\equiv0\pmod4,\\
 \tfrac14,&n\equiv1\pmod4,\\
 \tfrac12,&n\equiv2\pmod4,\\
 \tfrac34,&n\equiv3\pmod4.
 \end{cases}
\end{equation}
The same periodic correction is
\[
 \frac{5+(-1)^n+(1-i)(-i)^n+(1+i)i^n}{8},\qquad i^2=-1.
\]
Substitution yields exactly the complex-valued-looking expression printed
in OEIS, whose total value is real and integral. This is an algebraic
identity for $q$, independent of the unresolved equality $a=q$.

For all $n\ge0$, another convenient form is
\[
 q(n)=\frac{n(n+1)}2-\left\lfloor\frac{n-1}{4}\right\rfloor
                         +\ind{1\le n\le7}.
\]
Here $\lfloor-1/4\rfloor=-1$ correctly gives $q(0)=1$.
The first terms are
\[
 1,2,4,7,11,15,21,28,35,43,53,64,76,88,102,117,133,\ldots.
\]

\subsection{Why the recurrence starts at fourteen}
Set
\begin{align*}
 D(z)&=(1-z)^2(1-z^4)=1-2z+z^2-z^4+2z^5-z^6,\\
 N(z)&=1+z^2+z^3+z^6-z^8+z^9+z^{12}-z^{13}.
\end{align*}
The identity $D(z)Q(z)=N(z)$ implies, for $n\ge6$,
\begin{equation}\label{eq:inhomogeneous}
 q(n)-2q(n-1)+q(n-2)-q(n-4)+2q(n-5)-q(n-6)
 =\coef{z}{n}N(z).
\end{equation}
The right side vanishes for every $n\ge14$, but at $n=13$ it equals $-1$.
Thus the least starting index from which the homogeneous recurrence holds
for \emph{every} subsequent index of $q$ is $14$.
For $7\le n\le13$, its only nonzero residuals are
\begin{center}
\begin{tabular}{c|rrrr}
$n$&8&9&12&13\\\hline
residual&$-1$&1&1&$-1$
\end{tabular}
\end{center}
This proves the range correction asserted in Section 1.
If the original derivative sequence satisfies the corrected recurrence
for all $n\ge14$, its directly verified initial values $a(0),\ldots,a(13)$
force $a=q$ by induction. Consequently the corrected recurrence for $a$
is itself equivalent to the main conjecture, once those initial values
are supplied; it is not an independent proved consequence for arbitrary
orders.

The second forward differences $q(n+2)-2q(n+1)+q(n)$ begin
\[
 1,1,1,0,2,1,0,1,\underbrace{2,1,1,0,2,1,1,0,\ldots}_{n\ge8}.
\]
Again, this is proved for $q$; transferring the entire pattern to $a$
requires the missing nonvanishing statement.

\section{Reproducible verification}\label{sec:computation}

\subsection{What was actually checked}
All reported checks use exact arithmetic. No floating-point threshold was
used to decide that a coefficient was zero.

The C++ implementation with GMP computes the recurrence \eqref{eq:brec}
through row $3000$. It tests all $4{,}501{,}500$ entries with
$1\le r\le m\le3000$. Exactly $750$ are zero: $749$ symmetry zeros and
$(8,5)$. There are no other zeros in that range, and therefore
\begin{equation}\label{eq:3000}
 a(n)=q(n)\quad(0\le n\le3000),\qquad
 a(3000)=4{,}500{,}751.
\end{equation}
These are finite computational results, not an induction proving all $n$.

A separate Python standard-library implementation computes the triangle
through row $1500$ and agrees with the C++ term counts at all $1501$
indices $0,\ldots,1500$. Further independent checks include the complete
polynomials from the differential recurrence \eqref{eq:Prec} through order
$60$, all $496$ shifted-product coefficients \eqref{eq:falling} with
$0\le r\le m\le30$, and $11{,}865$ instances of the prime-adic valuation
formula. All $54$ values displayed directly in the OEIS entry agree with
the exact computation. The separately linked OEIS table through $500$ was
not used as an independent data source for this report.

The sixteen diagonal certificates were reconstructed and checked with a
third, rational-polynomial computation. In addition to showing no unexpected
zeros in a finite triangular region, these certificates prove there are no
unexpected zeros on those fixed diagonals for any row index.

\subsection{Algorithm and complexity}
Only two preceding rows are required for \eqref{eq:brec}. A schematic
implementation is
\begin{quote}\ttfamily\small
previousprevious = [1]\\
previous = [0, 1]\\
count = 2\\
for m = 2, ..., N:\\
\hspace*{1em}current[0] = 0\\
\hspace*{1em}for r = 1, ..., m:\\
\hspace*{2em}current[r] = (r-m+1)*previous[r]\\
\hspace*{3em}+ previous[r-1] + (m-1)*previousprevious[r-1]\\
\hspace*{1em}count += number of nonzero current[r] for 1 <= r <= m\\
\hspace*{1em}rotate the three rows
\end{quote}
Indices beyond the end of a row mean zero. The supplied implementations
handle these boundaries explicitly.

The computation uses $O(N^2)$ arithmetic operations and stores $O(N)$
integers. These are not bit-complexity claims: the integers grow with $N$.
From the product formula,
\[
 |b(m,r)|\le\prod_{h=0}^{m-1}(1+|r-h|)\le(m+1)^m,
\]
so a coefficient has $O(m\log m)$ bits. The rolling-row implementation thus
has an $O(N^2\log N)$-bit storage bound. A crude schoolbook arithmetic bound
for the running time is $O(N^3\log^2 N)$ bit operations.

\subsection{Files and commands}
The archive contains \texttt{article.tex}, the compiled \texttt{article.pdf},
a build script, both triangle verifiers, the rational certificate verifier,
the complete count and zero-position tables through $3000$, and the exact
verification reports. From the extracted directory, run
\begin{quote}\ttfamily\small
python3 verify.py --max-n 1500 --data-dir data\\
python3 verify\_diagonals.py --output-dir data --check-only
\end{quote}
The C++ implementation requires a C++17 compiler and GMP development files:
\begin{quote}\ttfamily\small
c++ -O3 -std=c++17 verify\_gmp.cpp -lgmpxx -lgmp -o verify\_gmp\\
./verify\_gmp 3000 data
\end{quote}
The Python programs use only the standard library. The diagonal checker
constructs the coefficients afresh rather than trusting the saved residue
tables. Its proof obligations are elementary rational polynomial identities
and finite modular evaluations; it does not require a general-purpose
computer algebra system.

\section{What remains to finish the requested proof}

The derivation of the candidate rational function is complete, but the
identification of that function with the derivative count is conditional.
The exact outstanding statement is
\[
 \coef{u}{r}\prod_{h=0}^{m-1}(u+r-h)\ne0
\]
for all $1\le r\le m$ outside the explicitly proved zero set $\ZZ_0$.
The preceding results settle this for $r\le3$, for $m-r\le16$, and whenever
$m-r+1$ is prime, in addition to all rows through $3000$ by exact computation.
They do not cover every remaining pair.

The diagonal-polynomial formulation gives a particularly precise route to
a complete proof. For every offset $d$, it would suffice to show that the
residual polynomial $R_d$ in \eqref{eq:Rd} has no integer root $m>d$.
Irreducibility would be stronger than necessary. A uniform family of
modular nonvanishing certificates would also suffice, but the fact that
such certificates work for $d\le16$ does not establish that they work for
all $d$.

For the asymptotic conjecture alone, an exclusion of every additional zero
is unnecessary: by \eqref{eq:asympeq}, it suffices to bound their cumulative
number by $o(n^2)$. The present root-counting argument proves quadratic
order but not this sharper estimate.

\medskip
\noindent\textbf{Conclusion.}
The printed recurrence with range $n>6$ is disproved, and its correct
range for the proposed sequence is identified. The principal closed formula
is supported by an exact structural reduction, several infinite
nonvanishing theorems, and extensive reproducible verification. A complete
all-order proof of that formula is \emph{not} supplied: the missing converse
in Conjecture~\ref{conj:zeros} is left explicit rather than being replaced
by an inference from numerical agreement.

\clearpage
\begin{thebibliography}{9}
\bibitem{oeis293239}
V. Reshetnikov and contributors,
\emph{A293239: Number of terms in the fully expanded $n$th derivative of $x^x$},
The On-Line Encyclopedia of Integer Sequences, 2017.
Conjectured formulas attributed in the entry to C. Barker and V. Kote\v{s}ovec.
Consulted 20 September 2026.
\url{https://oeis.org/A293239}.

\bibitem{oeis8296}
N. J. A. Sloane and contributors,
\emph{A008296: Triangle of Lehmer--Comtet numbers of the first kind},
The On-Line Encyclopedia of Integer Sequences.
Consulted 20 September 2026.
\url{https://oeis.org/A008296}.

\bibitem{comtet}
L. Comtet,
\emph{Advanced Combinatorics: The Art of Finite and Infinite Expansions},
D. Reidel, Dordrecht, 1974, p. 139.

\bibitem{lehmer}
D. H. Lehmer,
\emph{Numbers associated with Stirling numbers and $x^x$},
Rocky Mountain Journal of Mathematics \textbf{15}(2) (1985), 461--479.
\url{https://projecteuclid.org/euclid.rmjm/1250127221}.

\bibitem{gould}
H. W. Gould,
\emph{A set of polynomials associated with the higher derivatives of $y=x^x$},
Rocky Mountain Journal of Mathematics \textbf{26}(2) (1996), 615--625.
\url{https://doi.org/10.1216/rmjm/1181072076}.

\bibitem{jeong}
S. Jeong,
\emph{On a question of Lehmer concerning the Comtet numbers},
arXiv:2607.26613v1, 2026.
In particular, Proposition 3.10(iii) records the symmetry zero family.
The recurrence question answered there concerns the numbers of the
\emph{second} kind, not the complete first-kind zero classification used here.
\url{https://arxiv.org/abs/2607.26613}.
\end{thebibliography}
\end{document}
